\section{Introduction}
\label{sec:intro}

This blueprint accompanies the
\href{https://github.com/bryangingechen/CombinatorialRigidity}{\texttt{CombinatorialRigidity}}
Lean~4 / mathlib4 formalization. The goal is a proof of \textbf{Laman's
theorem}~\cite{laman1970}: for $n \ge 2$, a graph $G$ on $n$ vertices is generically
rigid in the plane if and only if it contains a spanning subgraph $H$ with
$|E(H)| = 2n - 3$ such that every subgraph $H[S]$ on $k \ge 2$ vertices
satisfies $|E(H[S])| \le 2k - 3$. Such an $H$ is called a \emph{Laman
graph}, or \emph{minimally rigid graph in the plane}.

The Lean development follows the \emph{Henneberg-construction} route: we
characterize Laman graphs inductively via two combinatorial moves (Type~I
adds a vertex of degree~2; Type~II splits an edge), prove the moves
preserve generic rigidity, and combine these facts to deduce the ``Laman
implies generically rigid'' direction. This $\Leftarrow$ direction is the
classical strategy of Tay and Whiteley~\cite{tayWhiteley1985}, in the
modern presentation of Jord\'an~\cite{jordan2016}~\S 2.2. The converse
direction goes through the rigidity matroid. For background on
combinatorial rigidity, see
\cite{graverServatiusServatius1993,jordan2016}.

\subsection*{Phase plan}

The Lean code is organized into seven phases, tracked in
\href{https://github.com/bryangingechen/CombinatorialRigidity/blob/master/ROADMAP.md}{\texttt{ROADMAP.md}}.
All seven phases are complete: the first six phases formalize
Laman's theorem, and Phase~7 packages the planar rigidity matroid
as a mathlib \texttt{Matroid} via the Lov\'asz--Yemini matroid
identification.

\begin{enumerate}
  \item \textbf{Sparsity.} Basic $(k,\ell)$-sparsity and $(k,\ell)$-tightness
        for arbitrary $k, \ell$. (Files: \texttt{EdgesIn.lean},
        \texttt{Sparsity.lean}.)
  \item \textbf{Laman graphs.} The $(2,3)$-tight specialization, the
        $K_2$ base case, and degree facts feeding the Henneberg induction.
        (File: \texttt{Laman.lean}.)
  \item \textbf{Henneberg moves.} Type~I and Type~II moves, Laman-preservation
        in both directions, and the existence of an inverse move at every
        non-base Laman graph. (File: \texttt{Henneberg.lean}.)
  \item \textbf{Frameworks.} Frameworks, the rigidity map, infinitesimal
        rigidity, generic rigidity. (File: \texttt{Framework.lean}.)
  \item \textbf{Laman's theorem ($\Leftarrow$).} Per-move rigidity
        preservation and the induction. (Files: \texttt{HennebergRigidity.lean},
        \texttt{LamanTheorem.lean}.)
  \item \textbf{Laman's theorem ($\Rightarrow$).} The companion direction
        via the planar rigidity matroid (Lov\'asz--Yemini's easy
        identification of row-independence with $(2,3)$-sparsity).
        (Files: \texttt{RigidityMatroid.lean}, \texttt{LamanTheorem.lean}.)
  \item \textbf{Lov\'asz--Yemini matroid identification.} The harder
        converse: $(2,3)$-sparse $\Rightarrow$ row-independent at some
        placement; together with Phase~6 this gives the
        Lov\'asz--Yemini~\cite{lovaszYemini1982} row-LI $\Leftrightarrow$
        $(2,3)$-sparse identification. Packaged as a mathlib
        \texttt{Matroid} via the general $(k, \ell)$-count matroid in
        the matroidal regime $\ell < 2k$
        (Whiteley~\cite{whiteley1996}, Lee--Streinu~\cite{leeStreinu2008}).
        (Files: \texttt{CountMatroid.lean}, \texttt{MatroidIdentification.lean}.)
\end{enumerate}

The \emph{linear-matroid} framing of the planar rigidity matroid (the
matroid of the rigidity-row function at a generic placement, with
Lov\'asz--Yemini as a matroid isomorphism with the
$(2, 3)$-count matroid) is left to a follow-up phase.

\subsection*{Reading this blueprint}

Statements that are formalized in Lean carry a green tick and a
\texttt{\textbackslash lean\{\}} pointer linking back to the API
documentation. The \emph{dependency graph view} (linked from the
navigation bar) is the best place to see at a glance how the chapters
fit together; nodes on the main Laman line and through the
Phase~7 matroid identification are all green.
